\documentclass{article}
\usepackage{enumitem}

\title{Medical Case Study: Atypical Presentation of Lyme Disease}
\author{Internal Medicine, Community Hospital}

\begin{document}
\maketitle

\section*{Patient Presentation}
A 34-year-old female presented with a two-week history of intermittent fevers, diffuse joint pain, and fatigue. No erythema migrans was observed.

\subsection*{Past medical history}
\begin{itemize}[noitemsep]
  \item No significant past medical history
  \item No recent travel
  \item Resides in Lyme-endemic region
\end{itemize}

\subsection*{Examination}
Vital signs within normal limits. Physical exam revealed mild synovitis in right knee. Skin exam unremarkable.

\subsection*{Investigations}
Initial workup:
\begin{itemize}[noitemsep]
  \item CBC, CMP: unremarkable
  \item ESR: 32 mm/hr (elevated)
  \item CRP: 18 mg/L (elevated)
  \item Lyme ELISA: positive
  \item Western blot (IgG): positive
\end{itemize}

\section*{Diagnosis}
Late-stage disseminated Lyme disease with monoarticular arthritis.

\section*{Treatment}
28-day course of doxycycline 100mg BID.

\section*{Follow-up}
At 8-week follow-up, arthritis resolved. ESR and CRP normalised.

\section*{Discussion}
This case highlights that Lyme disease can present without the classic erythema migrans rash. Clinicians in endemic regions should maintain a broad differential.

\end{document}
